\section{特定業種退職金共済}

\begin{fillblank}[年金1 2018 問題1(5)\textcircled{4}]
  特定業種退職金共済に加入できる企業は、建設業、清酒製造業、（　　A　　）を営む全ての事
業主である。
\end{fillblank}

\begin{answerbox}
  \begin{enumerate}[A:]
    \item 林業
  \end{enumerate}
\end{answerbox}

\begin{fillblank}[年金1 2023 問題1(6)(ア)]
（　　A　　）制度は、 （　　B　　）・清酒製造業・林業を営む企業向けに設けられた社外積立の
退職給付制度であり、 （　　C　　）制度とは別に、 （　　C　　）法（昭和34 年法律第160 号）
に基づいて独立行政法人（　　D　　）が運営を行っている。
\end{fillblank}

\begin{answerbox}
  \begin{enumerate}[A:]
    \item 特定業種退職金共済
    \item 建設業
    \item 中小企業退職金共済
    \item 勤労者退職金共済機構
  \end{enumerate}
\end{answerbox}

\begin{fillblank}[年金1 2021 問題1(6)\textcircled{3}（改訂）]
  特定業種退職金共済の掛金の日額は以下の通りである。

  \begin{itemize}
    \item 建設業：（　　A　　）円
    \item 清酒製造業：（　　B　　）円
    \item 林業：（　　C　　）円
  \end{itemize}
\end{fillblank}

\begin{answerbox}
  \begin{enumerate}[A:]
    \item 320
    \item 300
    \item 470
  \end{enumerate}
\end{answerbox}

\begin{fillblank}[年金1 2024 問題1(5)ア]
  特定業種退職金共済の退職金は、共済証紙の納付月数が（　　A　　）月（建設業退職金共済およ
び死亡の場合は（　　B　　）月）以上である者が、他の業種に従事したときや５５歳以上で退職
したときなどに該当したときに、納付月数に応じて一時金で支給される。
\end{fillblank}

\begin{answerbox}
  \begin{enumerate}[A:]
    \item 24
    \item 12
  \end{enumerate}
\end{answerbox}

\newpage